% Luc Destombe --- licence CC BY-NC-SA 4.0
% https://creativecommons.org/licenses/by-nc-sa/4.0/deed.fr
% Fichier autonome : compiler avec pdflatex, deux fois (signets, nombre de pages).
\def\lyceeversion{1.0}
\documentclass[a4paper,10pt]{article}
\makeatletter
\newif\iflycee@niveaux
\lycee@niveauxtrue
\newif\iflycee@cm
\newif\iflycee@palatino

\RequirePackage[utf8]{inputenc}
\RequirePackage[T1]{fontenc}
\RequirePackage[french]{babel}
\RequirePackage{etoolbox}

\newcommand{\ShorthandsOff}{\@ifundefined{shorthandoff}{}{\shorthandoff{;:!?}}}
\newcommand{\ShorthandsOn}{\@ifundefined{shorthandon}{}{\shorthandon{;:!?}}}
\ShorthandsOff
\AtBeginDocument{\ShorthandsOn}
\AtBeginEnvironment{tikzpicture}{\ShorthandsOff}

\RequirePackage{geometry}
\RequirePackage{fancyhdr}
\RequirePackage{lastpage}
\RequirePackage{multicol}
\RequirePackage{array}
\RequirePackage{enumitem}
\RequirePackage{xcolor}
\RequirePackage{needspace}

\RequirePackage{amsmath,amssymb}
\iflycee@cm\else\RequirePackage{fourier}\fi
\DeclareMathAlphabet{\mathcal}{OMS}{cmsy}{m}{n}
\renewcommand{\ttdefault}{lmtt}
\begingroup\catcode`\_=\active \catcode`\^=\active
\global\let_\sb \global\let^\sp
\endgroup
\AtBeginDocument{\catcode`\_=12 \mathcode`\_="8000
  \catcode`\^=12 \mathcode`\^="8000 }
\newlength{\retraitcalculs}
\setlength{\retraitcalculs}{2em}

\RequirePackage{tikz}
\usetikzlibrary{arrows.meta,calc,babel,shapes.misc}
\RequirePackage[most]{tcolorbox}
\tcbuselibrary{skins,breakable}

\definecolor{coulDef}{HTML}{1F4E79}
\definecolor{coulProp}{HTML}{7030A0}
\definecolor{coulRem}{HTML}{C55A11}
\definecolor{coulEx}{HTML}{2E7D32}
\definecolor{coulExo}{HTML}{B03A2E}
\definecolor{coulPart}{HTML}{1F4E79}
\definecolor{coulMeth}{HTML}{1B6E4F}
\definecolor{coulCourbe}{HTML}{0B5ED7}
\definecolor{coulCourbeB}{HTML}{C00000}
\definecolor{coulCourbeC}{HTML}{2E7D32}
\definecolor{coulGrille}{HTML}{8C8C8C}
\definecolor{coulRep}{HTML}{1B6E4F}
\definecolor{coulBar}{HTML}{2E7D32}

\newcommand{\R}{\mathbb{R}}
\newcommand{\Rs}{\mathbb{R}^{*}}
\newcommand{\N}{\mathbb{N}}
\newcommand{\Ns}{\mathbb{N}^{*}}
\newcommand{\Z}{\mathbb{Z}}
\newcommand{\Q}{\mathbb{Q}}
\renewcommand{\le}{\leqslant}
\renewcommand{\ge}{\geqslant}
\renewcommand{\approx}{\simeq}
\newcommand{\vect}[1]{\overrightarrow{#1}}
\newcommand{\norme}[1]{\left\lVert #1 \right\rVert}
\newcommand{\intervalle}[2]{\left[#1\,;#2\right]}
\RequirePackage{cancel}
\renewcommand{\CancelColor}{\color{coulExo}}
\newcommand{\fois}[1]{\times{\color{coulExo}#1}}
\newcommand{\barre}[1]{\cancel{#1}}
\newlength{\lycee@eqh}\newlength{\lycee@eqd}
\newsavebox{\lycee@eqa}\newsavebox{\lycee@eqb}\newsavebox{\lycee@eqc}
\newcommand{\lycee@eqcol}[1]{\hbox{$\renewcommand{\arraystretch}{1.35}%
  \begin{array}[t]{@{}>{\displaystyle}r@{}>{\displaystyle{}}l@{}}#1\end{array}$}}
\newcommand{\lycee@eqmesure}[1]{%
  \ifdim\ht#1>\lycee@eqh\setlength{\lycee@eqh}{\ht#1}\fi
  \ifdim\dp#1>\lycee@eqd\setlength{\lycee@eqd}{\dp#1}\fi}
\newcommand{\lycee@eqtrait}{\hspace{1.2em}ou\hspace{1.2em}}
\newcommand{\lycee@eqiff}{$\iff$\hspace{0.8em}}
\newcommand{\eqproduit}[2]{\par\smallskip\noindent C'est une équation produit :
  $A\times B=0 \iff A=0$ ou $B=0$.\par\nopagebreak\smallskip
  \begingroup\setlength{\lycee@eqh}{0pt}\setlength{\lycee@eqd}{0pt}%
  \setbox\lycee@eqa=\lycee@eqcol{#1}\setbox\lycee@eqb=\lycee@eqcol{#2}%
  \lycee@eqmesure\lycee@eqa\lycee@eqmesure\lycee@eqb%
  \noindent\hspace*{\retraitcalculs}\lycee@eqiff\usebox\lycee@eqa\lycee@eqtrait\usebox\lycee@eqb\endgroup\par\smallskip}
\newcommand{\eqproduittrois}[3]{\par\smallskip\noindent C'est une équation produit :
  $A\times B\times C=0 \iff A=0$ ou $B=0$ ou $C=0$.\par\nopagebreak\smallskip
  \begingroup\setlength{\lycee@eqh}{0pt}\setlength{\lycee@eqd}{0pt}%
  \setbox\lycee@eqa=\lycee@eqcol{#1}\setbox\lycee@eqb=\lycee@eqcol{#2}\setbox\lycee@eqc=\lycee@eqcol{#3}%
  \lycee@eqmesure\lycee@eqa\lycee@eqmesure\lycee@eqb\lycee@eqmesure\lycee@eqc%
  \noindent\hspace*{\retraitcalculs}\lycee@eqiff\usebox\lycee@eqa\lycee@eqtrait\usebox\lycee@eqb\lycee@eqtrait\usebox\lycee@eqc\endgroup
  \par\smallskip}

\tikzset{grille/.style={coulGrille,line width=0.4pt}}
\newcommand{\quadrillage}[4]{\draw[grille,step=1] (#1,#2) grid (#3,#4);}
\tikzset{
  croix/.style={cross out,draw,line width=0.6pt,inner sep=0pt,minimum size=4pt},
  croixdroite/.style={croix,rotate=45,line width=0.8pt},
}
\newcommand{\pt}[3]{\path (#1) node[croix]{} node[#2,font=\small,inner sep=2.5pt]{$#3$};}

\newcommand{\lycee@repere}[4]{%
  \draw[grille,step=1] (#1,#2) grid (#3,#4);
  \draw[-{Stealth[length=2mm]},thick] (#1,0)--({#3+0.4},0) node[below left=-1pt]{$x$};
  \draw[-{Stealth[length=2mm]},thick] (0,#2)--(0,{#4+0.4}) node[below left=-1pt]{$y$};
  \node[below left,font=\scriptsize,inner sep=1.5pt] at (0,0) {$0$};}
\newcommand{\lycee@gradx}[1]{\foreach \k/\l in {#1}{%
  \draw (\k,0.07)--(\k,-0.07) node[below,font=\scriptsize,inner sep=1.5pt]{$\l$};}}
\newcommand{\lycee@grady}[1]{\foreach \k/\l in {#1}{%
  \draw (0.07,\k)--(-0.07,\k) node[left,font=\scriptsize,inner sep=1.5pt]{$\l$};}}
\newcommand{\lycee@intff}[2]{\left[#1\,;#2\right]}
\newcommand{\lycee@intfo}[2]{\left[#1\,;#2\right[}
\newcommand{\lycee@intof}[2]{\left]#1\,;#2\right]}
\newcommand{\lycee@intoo}[2]{\left]#1\,;#2\right[}
\newcommand{\lycee@ens}[1]{\left\{#1\right\}}
\AtBeginDocument{%
  \providecommand{\repere}{\lycee@repere}%
  \providecommand{\gradx}{\lycee@gradx}%
  \providecommand{\grady}{\lycee@grady}%
  \providecommand{\Cf}{\mathcal{C}\sb{f}}%
  \providecommand{\Cg}{\mathcal{C}\sb{g}}%
  \providecommand{\intff}{\lycee@intff}%
  \providecommand{\intfo}{\lycee@intfo}%
  \providecommand{\intof}{\lycee@intof}%
  \providecommand{\intoo}{\lycee@intoo}%
  \providecommand{\ens}{\lycee@ens}}

\newlist{questions}{enumerate}{3}
\setlist[questions,1]{label=\arabic*\textdegree),leftmargin=*,itemsep=2pt,topsep=3pt}
\setlist[questions,2]{label=\alph*),leftmargin=*,itemsep=1pt,topsep=2pt}
\setlist[questions,3]{label=\roman*),leftmargin=*,itemsep=1pt,topsep=1pt}
\setlist[itemize]{label=\textbullet,leftmargin=*,itemsep=2pt,topsep=3pt}

\newcommand{\besoinplace}[1]{\par\penalty\@M\begingroup
  \dimen@=#1\relax\dimen@ii=\pagegoal\advance\dimen@ii-\pagetotal
  \ifdim\dimen@>\dimen@ii\ifdim\dimen@ii>\z@\vfil\fi\break\fi\endgroup}

\newcommand{\lycee@pied}{\thepage/\pageref{LastPage}}
\newcommand{\entete}[2]{%
  \pagestyle{fancy}\fancyhf{}%
  \fancyhead[L]{\footnotesize\color{coulPart}\textbf{#1}}%
  \fancyhead[R]{\footnotesize\color{coulPart}\textbf{#2}}%
  \fancyfoot[C]{\footnotesize\lycee@pied}%
  \renewcommand{\headrulewidth}{0.4pt}%
  \renewcommand{\headrule}{\hbox to\headwidth{%
    \color{coulPart!40}\leaders\hrule height \headrulewidth\hfill}}}

\RequirePackage{ccicons}
\newcommand{\lycee@licence}{{\scriptsize\color{gray}%
  \copyright\ Luc Destombe --- \ccbyncsa\ CC BY-NC-SA 4.0}}
\newif\iflycee@licence \lycee@licencetrue
\newcommand{\sanslicence}{\lycee@licencefalse}
\AtBeginDocument{\iflycee@licence\AddToHookNext{shipout/foreground}{%
  \put(0,0){\rlap{\kern\dimexpr1in+\hoffset+\oddsidemargin+\textwidth\relax
    \llap{\raisebox{-\dimexpr1in+\voffset+\topmargin+\headheight+\headsep
      +\textheight+\footskip\relax}{\lycee@licence}}}}}\fi}

\newcommand{\lycee@site}{faire-des-maths.fr}
\newcommand{\lycee@version}{{\scriptsize\color{gray}\lycee@site
  \ifdefined\lyceeversion\ --- Version~\lyceeversion\fi}}
\AtBeginDocument{\iflycee@licence\AddToHookNext{shipout/foreground}{%
  \put(0,0){\rlap{\kern\dimexpr1in+\hoffset+\oddsidemargin\relax
    \raisebox{-\dimexpr1in+\voffset+\topmargin+\headheight+\headsep
      +\textheight+\footskip\relax}{\lycee@version}}}}\fi}

\newcommand{\formulesserrees}{%
  \setlength{\abovedisplayskip}{3pt}\setlength{\belowdisplayskip}{3pt}%
  \setlength{\abovedisplayshortskip}{2pt}\setlength{\belowdisplayshortskip}{3pt}}

\setlength{\parindent}{0pt}

\RequirePackage[implicit=false,hidelinks,pdfencoding=auto,psdextra,bookmarksopen,
  bookmarksopenlevel=1,pdfcreator={},pdfproducer={}]{hyperref}
\RequirePackage{bookmark}
\hypersetup{pdfauthor={Luc Destombe},pdfkeywords={CC BY-NC-SA 4.0}}
\newcounter{lycee@signet}
\newcommand{\signet}[3][]{\stepcounter{lycee@signet}%
  \edef\lycee@dest{\if\relax\detokenize{#1}\relax
    signet.\arabic{lycee@signet}\else#1\fi}%
  \hypertarget{\lycee@dest}{}\bookmark[level=#2,dest=\lycee@dest]{#3}}
\pdfstringdefDisableCommands{%
  \def\R{R}\def\N{N}\def\Z{Z}\def\Q{Q}\def\vect#1{#1}\def\Cf{Cf}\def\Cg{Cg}\def\fois#1{×#1}%
  \def\barre#1{#1}\def\intervalle#1#2{[#1;#2]}\def\intff#1#2{[#1;#2]}%
  \def\textsuperscript#1{#1}\def\dfrac#1#2{#1/#2}\def\frac#1#2{#1/#2}%
  \def\vec#1{#1⃗}}%
\RequirePackage[official]{eurosym}

\geometry{top=0.8cm,bottom=1.3cm,left=1cm,right=1cm,footskip=0.6cm}

\pagestyle{fancy}
\fancyhf{}
\renewcommand{\headrulewidth}{0pt}
\fancyfoot[C]{\footnotesize Page \thepage{} / \pageref{LastPage}}

\newcommand{\titrefiche}[2]{%
  \begin{center}
    \Large\textbf{#1}\\[2pt]
    \large\textbf{#2}\\[2pt]
  \end{center}
  \vspace{0.10cm}}

\setlength{\columnseprule}{0.4pt}
\setlength{\columnsep}{15pt}
\renewcommand{\columnseprulecolor}{\color{black!45}}
\setlength{\multicolsep}{2pt}

\newcounter{partiefiche}
\renewcommand{\thepartiefiche}{\Roman{partiefiche}}
\newif\ifapresbandeau
\newcommand{\lycee@bandeau}[1]{%
  \par\penalty-600
  \vspace{0.08cm}%
  \noindent\colorbox{coulPart}{%
    \parbox{\dimexpr\linewidth-2\fboxsep\relax}{%
      \color{white}\bfseries\raggedright\strut#1\strut}}%
  \nopagebreak\par\nopagebreak
  \vspace{0.06cm}\nopagebreak
  \apresbandeautrue}
\newcommand{\partiefiche}[1]{%
  \refstepcounter{partiefiche}%
  \lycee@bandeau{\signet{1}{\thepartiefiche{} --- #1}\thepartiefiche{} --- #1}}
\newcommand{\partiefichesansnum}[1]{\lycee@bandeau{\signet{1}{#1}#1}}

\newcounter{exo}
\newlength{\espaceexo}\setlength{\espaceexo}{7pt}
\newlength{\espacetitre}\setlength{\espacetitre}{2pt}
\newcommand{\exo}[1][\relax]{%
  \par
  \ifapresbandeau\apresbandeaufalse\else\penalty-150\addvspace{\espaceexo}\fi
  \refstepcounter{exo}%
  \noindent\signet[exercice-\arabic{exo}]{2}{Exercice \arabic{exo}}%
  \textbf{\color{coulExo}\underline{Exercice \theexo}}%
  \ifx\relax#1\relax\else\textbf{ : #1}\fi
  \par\nopagebreak\vspace{\espacetitre}\nopagebreak
  \ignorespaces}

\setlist[enumerate]{nosep,leftmargin=*,itemsep=2pt,topsep=2pt}
\setlist[itemize]{nosep,leftmargin=*,topsep=2pt}
\newcommand{\choix}[3]{\par\hspace*{1em}%
  \makebox[0.3\linewidth][l]{a) #1}\makebox[0.3\linewidth][l]{b) #2}c) #3}
\newcommand{\deux}[2]{\par\hspace*{0.6em}%
  \makebox[0.48\linewidth][l]{#1}#2}
\newcommand{\trois}[3]{\par\hspace*{0.6em}%
  \makebox[0.32\linewidth][l]{#1}\makebox[0.32\linewidth][l]{#2}#3}

  \renewcommand{\exo}[1][\relax]{%
    \ifapresbandeau\apresbandeaufalse\else\par\penalty-150\fi
    \refstepcounter{exo}%
    \vspace{0.05cm}%
    \noindent\signet[exercice-\arabic{exo}]{2}{Exercice \arabic{exo}}%
    \textbf{\color{coulExo}\underline{Exercice \theexo}}%
    \ifx\relax#1\relax\else\textbf{ : #1}\fi%
    \par\nopagebreak
    \ignorespaces}
  \newcommand{\rappel}[1]{%
    \par\vspace{0.06cm}\nopagebreak
    \noindent\fcolorbox{black!35}{black!5}{%
      \parbox{\dimexpr\linewidth-2\fboxsep-2\fboxrule\relax}{%
        \small\textbf{Rappel.} #1}}%
    \par\vspace{0.06cm}\nopagebreak}
  \setlist[enumerate]{nosep,leftmargin=*}
  \setlist[itemize]{nosep,leftmargin=*}
  \setlength{\multicolsep}{3pt plus 1pt minus 1pt}
  \setlength{\premulticols}{10pt}
  \setlength{\postmulticols}{10pt}
  \newlist{expr}{enumerate}{1}
  \setlist[expr]{label={\Alph*\ $=$},nosep,leftmargin=*,itemsep=0pt}

\renewenvironment{center}{\par\addvspace{3pt}\centering}{\par\addvspace{3pt}}
\formulesserrees
\AtBeginDocument{\formulesserrees}
\clubpenalty=10000
\widowpenalty=10000
\displaywidowpenalty=10000
\brokenpenalty=10000
\setlength{\parskip}{0pt}

\RequirePackage{ragged2e}
\setlength{\RaggedRightRightskip}{0pt plus 1fil}
\AtBeginDocument{\RaggedRight\binoppenalty=10000 \relpenalty=10000 }
\g@addto@macro\@parboxrestore{\RaggedRight}
\makeatother
\usepackage{tkz-tab}
\geometry{top=0.8cm,bottom=1.3cm,left=1cm,right=1cm,footskip=0.7cm,headheight=12pt}

\newcommand{\e}{\mathrm{e}}

\newtcblisting{python}{%
  listing only,enhanced,breakable,
  colback=black!3,colframe=black!35,boxrule=0.5pt,arc=1pt,
  left=6pt,right=6pt,top=2pt,bottom=2pt,
  listing options={language=Python,basicstyle=\ttfamily\small,
    keywordstyle=\color{coulMeth}\bfseries,columns=fullflexible,
    keepspaces=true,showstringspaces=false,aboveskip=0pt,belowskip=0pt}}

\begin{document}
\titrefiche{Terminale spécialité --- Dérivation, convexité, continuité}{Fiche d'exercices du chapitre}

\begin{multicols}{2}

\partiefiche{Rappels sur la dérivation}

\exo[Lire des nombres dérivés et des tangentes]
On donne la courbe $\Cf$ d'une fonction $f$ définie sur $\intff{-1}{4}$, et ses
tangentes aux points d'abscisses $0$, $1$ et $3$ (en pointillés).
\begin{center}
\begin{tikzpicture}[x=0.55cm,y=0.55cm]
  \repere{-2}{-2}{5}{4}
  \gradx{-1,1,2,3,4} \grady{-1,1,2,3}
  \draw[coulCourbeB,dashed,line width=0.8pt] (-1.2,-1.4)--(1.2,3.4) (-0.5,2)--(2.5,2)
    (1.2,1.8)--(4.6,-1.6);
  \draw[coulCourbe,line width=1.1pt]
    (-1,-1) .. controls (-0.67,-0.33) and (-0.33,0.33) .. (0,1)
            .. controls (0.33,1.67) and (0.67,2) .. (1,2)
            .. controls (1.67,2) and (2.33,0.67) .. (3,0)
            .. controls (3.33,-0.33) and (3.67,-0.67) .. (4,-1);
  \path (0,1) node[croix]{} (1,2) node[croix]{} (3,0) node[croix]{};
  \node[coulCourbe,font=\footnotesize] at (-0.9,0.4) {$\Cf$};
\end{tikzpicture}
\end{center}
\begin{enumerate}
  \item Lire $f(0)$, $f(1)$ et $f(3)$.
  \item Lire $f'(0)$, $f'(1)$ et $f'(3)$.
  \item En déduire une équation de chacune des tangentes $(T_{0})$, $(T_{1})$ et $(T_{3})$.
\end{enumerate}

\exo[Calculer des dérivées]
Calculer la dérivée de chaque fonction, et l'écrire sous forme factorisée quand c'est
possible.
\begin{enumerate}[label=\alph*)]
  \item $f$ définie sur $\R$ par $f(x)=\left(x^{2}-3\right)\e^{x}$.
  \item $g$ définie sur $\intoo{-\frac{3}{2}}{+\infty}$ par $g(x)=\dfrac{4x-1}{2x+3}$.
  \item $h$ définie sur $\R$ par $h(x)=\dfrac{5}{x^{2}+1}$.
  \item $k$ définie sur $\R$ par $k(x)=x\e^{-3x}$.
  \item $m$ définie sur $\intoo{-1}{+\infty}$ par $m(x)=\dfrac{\e^{x}}{x+1}$.
\end{enumerate}

\exo[Tangentes à une courbe]
Soit la fonction $f$ définie sur $\R$ par $f(x)=(x+3)\e^{-x}$, de courbe $\Cf$.
\begin{enumerate}
  \item Montrer que, pour tout $x\in\R$, $f'(x)=(-x-2)\e^{-x}$.
  \item Déterminer une équation de la tangente $(T_{0})$ à $\Cf$ au point d'abscisse $0$.
  \item La tangente $(T_{0})$ passe-t-elle par le point $A(2\,;-1)$ ?
  \item En quel point la courbe $\Cf$ admet-elle une tangente horizontale ? Donner ses
        coordonnées exactes.
\end{enumerate}

\exo[Une étude de variations]
Soit la fonction $f$ définie sur $\R$ par $f(x)=(2x-5)\e^{x}$.
\begin{enumerate}
  \item Déterminer les limites de $f$ en $-\infty$ et en $+\infty$.
  \item Calculer $f'(x)$ pour tout $x\in\R$.
  \item Dresser le tableau de variations de $f$, avec la valeur exacte de son minimum.
\end{enumerate}

\partiefiche{Fonctions composées}

\exo[Composer, décomposer]
\begin{enumerate}
  \item Pour chaque couple de fonctions $u$ et $v$, exprimer $v\circ u(x)$ et
        $u\circ v(x)$ en fonction de $x$.
  \begin{enumerate}[label=\alph*)]
    \item $u$ et $v$ définies sur $\R$ par $u(x)=x+2$ et $v(x)=x^{2}$.
    \item $u$ et $v$ définies sur $\R$ par $u(x)=x^{2}$ et $v(x)=\e^{x}$.
  \end{enumerate}
  \item Soit les fonctions $u$ et $v$ définies par $u(x)=1-x$ sur $\R$ et $v(x)=\sqrt{x}$
        sur $\intfo{0}{+\infty}$. Pour quels nombres réels $x$ la composée $v\circ u$
        est-elle définie ? Exprimer $v\circ u(x)$.
  \item Écrire chaque fonction, définie sur $\R$, comme la composée de deux fonctions
        $u$ et $v$ à préciser.
  \begin{enumerate}[label=\alph*)]
    \item $f(x)=(5x-2)^{3}$
    \item $g(x)=\sqrt{x^{2}+1}$
    \item $h(x)=\e^{3x^{2}-x}$
  \end{enumerate}
\end{enumerate}

\exo[Dériver une exponentielle composée]
Calculer la dérivée de chaque fonction.
\begin{enumerate}[label=\alph*)]
  \item $f$ définie sur $\R$ par $f(x)=\e^{4x-1}$.
  \item $g$ définie sur $\R$ par $g(x)=\e^{x^{2}+1}$.
  \item $h$ définie sur $\R$ par $h(x)=3\e^{-x^{2}}$.
  \item $k$ définie sur $\intoo{0}{+\infty}$ par $k(x)=\e^{\frac{1}{x}}$.
  \item $m$ définie sur $\R$ par $m(x)=x\e^{x^{2}}$.
\end{enumerate}

\exo[Puissances et racines]
Calculer la dérivée de chaque fonction.
\begin{enumerate}[label=\alph*)]
  \item $f$ définie sur $\R$ par $f(x)=(2x+7)^{4}$.
  \item $g$ définie sur $\R$ par $g(x)=\left(x^{2}-x+3\right)^{3}$.
  \item $h$ définie sur $\intoo{\frac{1}{3}}{+\infty}$ par $h(x)=\dfrac{1}{(3x-1)^{2}}$.
  \item $k$ définie sur $\intoo{-\frac{1}{3}}{+\infty}$ par $k(x)=\sqrt{6x+2}$.
  \item $m$ définie sur $\R$ par $m(x)=\sqrt{x^{2}+4x+5}$ (vérifier d'abord que
        $x^{2}+4x+5>0$).
  \item $p$ définie sur $\R$ par $p(x)=\left(1+2\e^{x}\right)^{3}$.
\end{enumerate}

\exo[Tangentes à la courbe d'une racine]
Soit la fonction $f$ définie sur $\intfo{-\frac{1}{2}}{+\infty}$ par $f(x)=\sqrt{2x+1}$,
de courbe $\Cf$ ; on admet qu'elle est dérivable sur $\intoo{-\frac{1}{2}}{+\infty}$.
\begin{enumerate}
  \item Montrer que, pour tout $x>-\frac{1}{2}$, $f'(x)=\dfrac{1}{\sqrt{2x+1}}$.
  \item Déterminer une équation de la tangente $(T_{0})$ à $\Cf$ au point d'abscisse $0$,
        puis de la tangente $(T_{4})$ au point d'abscisse $4$.
  \item Existe-t-il un point de $\Cf$ où la tangente a pour coefficient directeur
        $\frac{1}{5}$ ? Si oui, donner ses coordonnées.
\end{enumerate}

\exo[Repérer les erreurs]
Un élève a écrit les dérivées suivantes, toutes fausses. Pour chacune, expliquer
l'erreur en une phrase, puis écrire la bonne dérivée.
\begin{enumerate}[label=\alph*)]
  \item $f$ définie sur $\R$ par $f(x)=\e^{x^{2}}$ : « $f'(x)=\e^{x^{2}}$ ».
  \item $g$ définie sur $\R$ par $g(x)=(3x+1)^{4}$ : « $g'(x)=4(3x+1)^{3}$ ».
  \item $h$ définie sur $\R$ par $h(x)=\sqrt{x^{2}+5}$ :
        « $h'(x)=\dfrac{1}{2\sqrt{x^{2}+5}}$ ».
\end{enumerate}

\partiefiche{Étude d'une fonction composée}

\exo[Une exponentielle de trinôme]
Soit la fonction $f$ définie sur $\R$ par $f(x)=\e^{-x^{2}+4x-3}$.
\begin{enumerate}
  \item Déterminer les limites de $f$ en $-\infty$ et en $+\infty$. Interpréter
        graphiquement.
  \item Calculer $f'(x)$ pour tout $x\in\R$.
  \item Dresser le tableau de variations de $f$.
  \item Résoudre l'équation $f(x)=1$.
\end{enumerate}

\exo[La racine d'un trinôme]
Soit la fonction $g$ définie par $g(x)=\sqrt{x^{2}-6x+10}$.
\begin{enumerate}
  \item Montrer que, pour tout $x\in\R$, $x^{2}-6x+10=(x-3)^{2}+1$. En déduire que $g$
        est définie et dérivable sur $\R$.
  \item Déterminer les limites de $g$ en $-\infty$ et en $+\infty$.
  \item Calculer $g'(x)$ pour tout $x\in\R$.
  \item Dresser le tableau de variations de $g$.
\end{enumerate}

\exo[La puissance d'une installation solaire]
Par une journée ensoleillée, la puissance produite par les panneaux solaires d'une
école, en kilowatts (kW), est modélisée par la fonction $P$ définie sur $\intff{6}{20}$
par
\begin{center}
$P(t)=8\e^{-\frac{(t-13)^{2}}{8}}$,
\end{center}
où $t$ est l'heure de la journée.
\begin{enumerate}
  \item Montrer que, pour tout $t\in\intff{6}{20}$,
        $P'(t)=-2(t-13)\e^{-\frac{(t-13)^{2}}{8}}$.
  \item Étudier les variations de $P$ sur $\intff{6}{20}$.
  \item À quelle heure la puissance est-elle maximale ? Quelle est cette puissance ?
  \item À l'aide de la calculatrice, déterminer à partir de quelle heure, à la minute
        près, la puissance dépasse $4$~kW.
\end{enumerate}

\partiefiche{Dérivée seconde}

\exo[Calculer des dérivées secondes]
Calculer la dérivée seconde de chaque fonction.
\begin{enumerate}[label=\alph*)]
  \item $f$ définie sur $\R$ par $f(x)=x^{4}-2x^{3}+x-5$.
  \item $g$ définie sur $\R$ par $g(x)=\e^{3x}+x^{2}$.
  \item $h$ définie sur $\R$ par $h(x)=(x+1)\e^{2x}$.
  \item $k$ définie sur $\intoo{0}{+\infty}$ par $k(x)=\dfrac{1}{x}$.
  \item $m$ définie sur $\R$ par $m(x)=x\e^{-x}$.
\end{enumerate}

\exo[Dérivée et dérivée seconde]
Soit la fonction $f$ définie sur $\R$ par $f(x)=\left(x^{2}+1\right)\e^{-x}$.
\begin{enumerate}
  \item Montrer que, pour tout $x\in\R$, $f'(x)=-(x-1)^{2}\e^{-x}$.
  \item En déduire le sens de variation de $f$ sur $\R$.
  \item Montrer que, pour tout $x\in\R$, $f''(x)=(x-1)(x-3)\e^{-x}$.
  \item Étudier les variations de $f'$ sur $\R$.
\end{enumerate}

\exo[Le démarrage d'un tramway]
Au démarrage, la distance parcourue par un tramway, en mètres, est modélisée par la
fonction $d$ définie sur $\intff{0}{10}$ par $d(t)=-0{,}2t^{3}+3t^{2}$, où $t$ est le
temps en secondes. Sa vitesse est $v(t)=d'(t)$, en m/s, et son accélération est
$a(t)=d''(t)$, en m/s$^{2}$.
\begin{enumerate}
  \item Calculer $v(t)$ et $a(t)$.
  \item À quel instant l'accélération est-elle nulle ?
  \item Étudier les variations de $v$ sur $\intff{0}{10}$. Quelle est la vitesse
        maximale du tramway ?
\end{enumerate}

\partiefiche{Fonctions convexes, fonctions concaves}

\exo[Lire la convexité sur trois courbes]
On donne la courbe $\Cf$ d'une fonction $f$, la courbe $\mathcal{C}_{g'}$ de la
\textbf{dérivée} d'une fonction $g$ et la courbe $\mathcal{C}_{h''}$ de la
\textbf{dérivée seconde} d'une fonction $h$.
\begin{center}
\begin{tikzpicture}[x=0.4cm,y=0.4cm]
  \repere{-1}{-1}{5}{5}
  \gradx{1,2,3,4} \grady{1,2,3,4}
  \draw[coulCourbe,line width=1.1pt]
    (0,0) .. controls (1,0) and (1.33,1.33) .. (2,2)
          .. controls (2.67,2.67) and (3,4) .. (4,4);
  \node[coulCourbe,font=\footnotesize] at (4.4,3.2) {$\Cf$};
\end{tikzpicture}\hfill
\begin{tikzpicture}[x=0.4cm,y=0.4cm]
  \repere{-3}{-4}{4}{2}
  \gradx{-2,-1,1,2,3} \grady{-3,-2,-1,1}
  \draw[coulCourbe,line width=1.1pt]
    (-2,-3) .. controls (-1.33,-1) and (-0.67,1) .. (0,1)
            .. controls (0.67,1) and (1.33,-0.11) .. (2,-1)
            .. controls (2.33,-1.44) and (2.67,-2.33) .. (3,-3);
  \node[coulCourbe,font=\footnotesize] at (3.3,-1.2) {$\mathcal{C}_{g'}$};
\end{tikzpicture}

\medskip
\begin{tikzpicture}[x=0.4cm,y=0.4cm]
  \repere{-4}{-3}{4}{3}
  \gradx{-3,-2,-1,1,2,3} \grady{-2,-1,1,2}
  \draw[coulCourbe,line width=1.1pt]
    (-3,-2) .. controls (-2.67,-1.17) and (-2.33,-0.44) .. (-2,0)
            .. controls (-1.33,0.89) and (-0.67,2) .. (0,2)
            .. controls (0.67,2) and (1.33,0.89) .. (2,0)
            .. controls (2.33,-0.44) and (2.67,-1.17) .. (3,-2);
  \node[coulCourbe,font=\footnotesize] at (1.6,2.6) {$\mathcal{C}_{h''}$};
\end{tikzpicture}
\end{center}
\begin{enumerate}
  \item Sur quels intervalles $f$ semble-t-elle convexe ? concave ?
  \item Sur quels intervalles $g$ est-elle convexe ? concave ? Justifier.
  \item Sur quels intervalles $h$ est-elle convexe ? concave ? Justifier.
\end{enumerate}

\exo[Du tableau de signes à la courbe]
Une fonction $f$, deux fois dérivable sur $\intff{-3}{5}$, a une dérivée seconde dont le
signe est donné par le tableau suivant.
\begin{center}
\begin{tikzpicture}[scale=0.8]
  \tkzTabInit[lgt=1.6,espcl=1.8]{$x$/0.8,$f''(x)$/0.8}{$-3$,$-1$,$3$,$5$}
  \tkzTabLine{,+,z,-,z,+,}
\end{tikzpicture}
\end{center}
\begin{enumerate}
  \item Déterminer la convexité de $f$ et les abscisses des points d'inflexion de sa
        courbe.
  \item Donner les variations de $f'$ sur $\intff{-3}{5}$.
  \item Est-il possible que $f'(-3)=2$ et $f'(-1)=0$ ? que $f'(-1)=1$ et $f'(3)=-2$ ?
  \item Sachant que $f(-1)=2$, $f(3)=-2$, $f'(-1)=1$ et $f'(3)=-2$, tracer dans un
        repère une courbe pouvant représenter $f$.
\end{enumerate}

\exo[Étudier la convexité]
Étudier la convexité de chaque fonction sur son ensemble de définition.
\begin{enumerate}[label=\alph*)]
  \item $f$ définie sur $\R$ par $f(x)=x^{3}-6x^{2}+2x-1$.
  \item $g$ définie sur $\R$ par $g(x)=x^{4}-6x^{2}+1$.
  \item $h$ définie sur $\R$ par $h(x)=\e^{-2x}+3x$.
  \item $k$ définie sur $\R$ par $k(x)=(2-x)\e^{x}$.
  \item $m$ définie sur $\intoo{0}{+\infty}$ par $m(x)=x+\dfrac{4}{x}$.
\end{enumerate}

\exo[D'après bac 2025 et 2026 : vrai ou faux ?]
Pour chaque affirmation, indiquer si elle est vraie ou fausse, en justifiant.
\begin{enumerate}
  \item On considère la fonction $f$ définie sur $\R$ par
        $f(x)=\left(-\dfrac{1}{2}x+3\right)^{5}$.

        \textbf{Affirmation 1 :} la fonction $f$ est convexe sur $\R$.
  \item On considère la fonction $g$ définie sur $\R$ par $g(x)=x\e^{x}$.

        \textbf{Affirmation 2 :} la fonction $g$ est convexe sur l'intervalle
        $\intfo{-2}{+\infty}$.
  \item On considère une fonction affine $h$ définie sur $\R$. On note $k$ la fonction
        définie sur $\R$ par $k(x)=x^{4}+x^{2}+h(x)$.

        \textbf{Affirmation 3 :} la fonction $k$ est convexe sur $\R$.
\end{enumerate}

\exo[D'après bac 2026, Amérique du Sud : QCM]
On donne ci-dessous la courbe représentative $\mathcal{C}_{f'}$ de la fonction dérivée
$f'$ d'une fonction $f$ définie et dérivable sur l'intervalle $\intff{-3}{4}$.
\begin{center}
\begin{tikzpicture}[x=0.45cm,y=0.45cm]
  \repere{-4}{-5}{5}{4}
  \gradx{-3,-2,-1,1,2,3,4} \grady{-4,-3,-2,-1,1,2,3}
  \draw[coulCourbe,line width=1.1pt]
    (-3,0) .. controls (-2.5,1.5) and (-2.25,3) .. (-2,3)
           .. controls (-1.75,3) and (-1.5,1.5) .. (-1,0)
           .. controls (-0.5,-1.5) and (-0.5,-4) .. (0,-4)
           .. controls (0.5,-4) and (1.5,-0.5) .. (2,0)
           .. controls (2.25,0.25) and (2.75,3) .. (3,3)
           .. controls (3.25,3) and (3.5,2) .. (4,1);
  \node[coulCourbe,font=\footnotesize] at (4.4,2.3) {$\mathcal{C}_{f'}$};
\end{tikzpicture}
\end{center}
Par lecture graphique, déterminer l'affirmation correcte pour la fonction $f$
(une seule est exacte), en justifiant.
\begin{enumerate}[label=\alph*)]
  \item La fonction $f$ est convexe sur l'intervalle $\intff{-2}{-1}$.
  \item La fonction $f$ est concave sur l'intervalle $\intff{-1}{2}$.
  \item La fonction $f$ est concave sur l'intervalle $\intff{-2}{0}$.
  \item La fonction $f$ admet exactement deux points d'inflexion sur l'intervalle
        $\intff{-3}{4}$.
\end{enumerate}

\partiefiche{Point d'inflexion}

\exo[Chercher les points d'inflexion]
Pour chaque fonction, définie sur $\R$, déterminer les coordonnées des éventuels points
d'inflexion de sa courbe.
\begin{enumerate}[label=\alph*)]
  \item $f(x)=x^{3}-3x^{2}+4$
  \item $g(x)=-x^{3}+6x^{2}-9x+2$
  \item $h(x)=x^{4}-4x^{3}$
  \item $k(x)=x^{4}+2x^{2}$
\end{enumerate}

\exo[Une vidéo de vulgarisation]
Le nombre de vues d'une vidéo de vulgarisation scientifique, en milliers, est modélisé
par la fonction $N$ définie sur $\intff{0}{10}$ par
\begin{center}
$N(t)=\dfrac{40}{1+\e^{4-t}}$,
\end{center}
où $t$ est le nombre de jours depuis sa mise en ligne.
\begin{enumerate}
  \item Montrer que, pour tout $t\in\intff{0}{10}$,
        $N'(t)=\dfrac{40\e^{4-t}}{\left(1+\e^{4-t}\right)^{2}}$, et en déduire le sens
        de variation de $N$.
  \item On admet que, pour tout $t\in\intff{0}{10}$,
        $N''(t)=\dfrac{40\e^{4-t}\left(\e^{4-t}-1\right)}{\left(1+\e^{4-t}\right)^{3}}$.
        Étudier la convexité de $N$ sur $\intff{0}{10}$.
  \item Montrer que la courbe de $N$ admet un point d'inflexion, dont on donnera les
        coordonnées.
  \item Interpréter ce résultat pour le nombre de vues de la vidéo.
\end{enumerate}

\exo[Deux points d'inflexion]
Soit la fonction $f$ définie sur $\R$ par $f(x)=x^{2}\e^{-x}$.
\begin{enumerate}
  \item Montrer que, pour tout $x\in\R$, $f'(x)=\left(2x-x^{2}\right)\e^{-x}$, puis que
        $f''(x)=\left(x^{2}-4x+2\right)\e^{-x}$.
  \item Étudier la convexité de $f$ sur $\R$.
  \item En déduire que la courbe de $f$ admet deux points d'inflexion, dont on donnera
        les abscisses exactes.
\end{enumerate}

\partiefiche{Convexité et inégalités}

\exo[Tangente et corde de l'exponentielle]
Soit la fonction $f$ définie sur $\R$ par $f(x)=\e^{x}$, de courbe $\Cf$.
\begin{enumerate}
  \item Justifier que $f$ est convexe sur $\R$.
  \item Déterminer une équation de la tangente $(T_{1})$ à $\Cf$ au point d'abscisse $1$.
        En déduire que, pour tout $x\in\R$, $\e^{x}\geq \e x$.
  \item Soit $A$ et $B$ les points de $\Cf$ d'abscisses $0$ et $1$. Déterminer une
        équation de la droite $(AB)$.
  \item En utilisant la convexité de $f$, en déduire que, pour tout
        $x\in\intff{0}{1}$, $\e^{x}\leq(\e-1)x+1$.
  \item En déduire, sans calculatrice, que $\sqrt{\e}\leq\dfrac{\e+1}{2}$.
\end{enumerate}

\exo[Au-dessus d'une tangente]
Soit la fonction $f$ définie sur $\R$ par $f(x)=x\e^{x}+x+1$, de courbe $\Cf$.
\begin{enumerate}
  \item Calculer $f'(x)$ et $f''(x)$ pour tout $x\in\R$.
  \item Étudier la convexité de $f$ sur $\R$.
  \item Déterminer une équation de la tangente $(T_{0})$ à $\Cf$ au point d'abscisse $0$.
  \item En déduire que, pour tout $x\geq -2$, $x\e^{x}+x+1\geq 2x+1$.
\end{enumerate}

\exo[Une approximation de racine carrée]
Soit la fonction $f$ définie sur $\intfo{-1}{+\infty}$ par $f(x)=\sqrt{1+x}$ ; on admet
qu'elle est deux fois dérivable sur $\intoo{-1}{+\infty}$.
\begin{enumerate}
  \item Montrer que, pour tout $x>-1$, $f''(x)=-\dfrac{1}{4(1+x)\sqrt{1+x}}$.
  \item En déduire la convexité de $f$ sur $\intoo{-1}{+\infty}$.
  \item Déterminer une équation de la tangente $(T_{0})$ à la courbe de $f$ au point
        d'abscisse $0$.
  \item En déduire que, pour tout $x>-1$, $\sqrt{1+x}\leq 1+\dfrac{x}{2}$.
  \item Sans calculatrice, en déduire un majorant de $\sqrt{1{,}02}$, puis le comparer à
        l'affichage de la calculatrice.
\end{enumerate}

\partiefiche{Continuité}

\exo[Le tarif d'un parking]
Dans un parking, on paie $2$~€ par heure commencée : jusqu'à $1$~h, on paie $2$~€ ; de
plus de $1$~h jusqu'à $2$~h, $4$~€, et ainsi de suite. On note $P(t)$ le prix payé, en
euros, pour une durée de $t$ heures, avec $t\in\intof{0}{4}$.
\begin{enumerate}
  \item Donner $P(0{,}5)$, $P(1)$, $P(1{,}5)$ et $P(2)$.
  \item Représenter la fonction $P$ dans un repère.
  \item Déterminer $\lim\limits_{t\to 1^{-}}P(t)$ et $\lim\limits_{t\to 1^{+}}P(t)$.
  \item La fonction $P$ est-elle continue en $1$ ? Justifier par une phrase complète.
  \item Sur quels intervalles la fonction $P$ est-elle continue ?
\end{enumerate}

\exo[Choisir un paramètre]
\begin{enumerate}
  \item Soit $a$ un nombre réel et $f$ la fonction définie sur $\R$ par
        $f(x)=x^{2}+1$ si $x<2$, et $f(x)=ax-3$ si $x\geq 2$. Déterminer $a$ pour que
        $f$ soit continue en $2$.
  \item Soit $b$ un nombre réel et $g$ la fonction définie sur $\R$ par
        $g(x)=\e^{x}$ si $x\leq 0$, et $g(x)=x+b$ si $x>0$. Déterminer $b$ pour que $g$
        soit continue sur $\R$.
\end{enumerate}

\exo[La fonction valeur absolue]
On rappelle que, pour tout $x\in\R$, $|x|=x$ si $x\geq 0$ et $|x|=-x$ si $x<0$. Soit la
fonction $f$ définie sur $\R$ par $f(x)=|x|$.
\begin{enumerate}
  \item Représenter la fonction $f$ dans un repère.
  \item Déterminer $\lim\limits_{x\to 0^{-}}f(x)$ et $\lim\limits_{x\to 0^{+}}f(x)$. La
        fonction $f$ est-elle continue en $0$ ?
  \item Pour $h\neq 0$, calculer le taux de variation $\dfrac{f(0+h)-f(0)}{h}$ selon le
        signe de $h$. La fonction $f$ est-elle dérivable en $0$ ?
  \item Que montre cet exemple sur le lien entre continuité et dérivabilité ?
\end{enumerate}

\exo[Image d'une suite convergente]
Soit la suite $(u_{n})$ définie pour tout $n\in\Ns$ par $u_{n}=3-\dfrac{1}{n}$.
\begin{enumerate}
  \item Déterminer la limite de la suite $(u_{n})$.
  \item En justifiant par la continuité, en déduire la limite de chacune des suites
        $(v_{n})$, $(w_{n})$ et $(t_{n})$ définies pour tout $n\in\Ns$ par
        $v_{n}=\e^{u_{n}}$, $w_{n}=\sqrt{u_{n}+1}$ et $t_{n}=u_{n}^{2}-2u_{n}$.
\end{enumerate}

\partiefiche{Théorème des valeurs intermédiaires}

\exo[Compter les solutions sur un tableau]
Une fonction $f$, continue sur $\intff{-3}{5}$, a le tableau de variations suivant.
\begin{center}
\begin{tikzpicture}[scale=0.8]
  \tkzTabInit[lgt=1.4,espcl=1.8]{$x$/0.8,$f(x)$/1.6}{$-3$,$1$,$3$,$5$}
  \tkzTabVar{-/$2$,+/$6$,-/$-2$,+/$4$}
\end{tikzpicture}
\end{center}
Déterminer, en justifiant, le nombre de solutions de chaque équation sur
$\intff{-3}{5}$.
\begin{enumerate}[label=\alph*)]
  \item $f(x)=0$
  \item $f(x)=5$
  \item $f(x)=3$
  \item $f(x)=-3$
\end{enumerate}

\exo[Une solution unique sur $\R$]
Soit la fonction $f$ définie sur $\R$ par $f(x)=x^{3}+x-3$.
\begin{enumerate}
  \item Étudier les variations de $f$ et ses limites en $-\infty$ et en $+\infty$.
  \item Montrer que l'équation $f(x)=0$ admet une unique solution $\alpha$ sur $\R$.
  \item Donner un encadrement de $\alpha$ d'amplitude $10^{-2}$.
  \item En déduire le signe de $f(x)$ selon les valeurs de $x$.
\end{enumerate}

\exo[Avec une exponentielle]
Soit la fonction $f$ définie sur $\R$ par $f(x)=\e^{x}+x-2$.
\begin{enumerate}
  \item Déterminer les limites de $f$ en $-\infty$ et en $+\infty$.
  \item Étudier le sens de variation de $f$ sur $\R$.
  \item Montrer que l'équation $f(x)=0$ admet une unique solution $\alpha$, et que
        $0<\alpha<1$.
  \item Donner un encadrement de $\alpha$ d'amplitude $10^{-2}$, puis le signe de $f(x)$
        selon les valeurs de $x$.
\end{enumerate}

\exo[D'après bac 2025, Amérique du Sud]
On se propose d'étudier la concentration dans le sang d'un médicament ingéré par une
personne pour la première fois. Soit $t$ le temps (en heures) écoulé depuis l'ingestion
de ce médicament. On admet que la concentration de ce médicament dans le sang, en gramme
par litre de sang, est modélisée par la fonction $f$ définie sur l'intervalle
$\intfo{0}{+\infty}$ par $f(t)=5t\e^{-t}$.
\begin{enumerate}
  \item Déterminer la limite de $f$ en $+\infty$. Interpréter ce résultat dans le
        contexte de l'exercice.
  \item Étudier les variations de $f$ sur l'intervalle $\intfo{0}{+\infty}$ puis dresser
        son tableau de variation complet.
  \item Démontrer qu'il existe deux nombres réels $t_{1}$ et $t_{2}$ tels que
        $f\left(t_{1}\right)=f\left(t_{2}\right)=1$. On donnera une valeur approchée à
        $10^{-2}$ des nombres réels $t_{1}$ et $t_{2}$.
  \item Pour une concentration du médicament supérieure ou égale à 1 gramme par litre de
        sang, il y a un risque de somnolence. Quelle est la durée en heures et minutes du
        risque de somnolence lors de la prise de ce médicament ?
\end{enumerate}

\exo[D'après bac 2025, Métropole]
On s'intéresse à la glycémie chez une personne venant de prendre un repas. La glycémie
en g.L$^{-1}$, en fonction du temps $t$, exprimé en heure, écoulé depuis la fin du
repas, est modélisée par la fonction $f$ définie sur $\intff{0}{6}$ par
$f(t)=(t+1)\e^{-0{,}4t}$.
\begin{enumerate}
  \item
  \begin{enumerate}[label=\alph*)]
    \item Montrer que, pour tout $t\in\intff{0}{6}$, $f'(t)=(-0{,}4t+0{,}6)\e^{-0{,}4t}$.
    \item Étudier les variations de $f$ sur $\intff{0}{6}$ puis dresser son tableau de
          variations sur cet intervalle.
  \end{enumerate}
  \item Une personne est en hypoglycémie lorsque sa glycémie est inférieure à
        $0{,}7$~g.L$^{-1}$.
  \begin{enumerate}[label=\alph*)]
    \item Démontrer que sur l'intervalle $\intff{0}{6}$ l'équation $f(t)=0{,}7$ admet une
          unique solution que l'on notera $\alpha$.
    \item Au bout de combien de temps après avoir pris son repas cette personne est-elle
          en hypoglycémie ? On exprimera ce temps à la minute près.
  \end{enumerate}
\end{enumerate}

\exo[D'après bac 2025, Asie : vrai ou faux ?]
Soit $f$ la fonction définie sur $\R$ par $f(x)=x\e^{-2x}$. On admet que $f$ est deux
fois dérivable sur $\R$. Pour chacune des affirmations suivantes, préciser si elle est
vraie ou fausse, puis justifier la réponse donnée.
\begin{enumerate}
  \item \textbf{Affirmation 1.} Pour tout nombre réel $x$, on a
        $f'(x)=(-2x+1)\e^{-2x}$.
  \item \textbf{Affirmation 2.} La fonction $f$ est convexe sur $\intof{-\infty}{1}$.
  \item \textbf{Affirmation 3.} L'équation $f(x)=-1$ admet une unique solution sur $\R$.
\end{enumerate}

\partiefiche{Méthode de dichotomie}

\exo[Dichotomie avec une boucle « pour »]
On cherche la solution de l'équation $\e^{x}=3$.
\begin{enumerate}
  \item Justifier que cette équation admet une unique solution $\alpha$ dans
        $\intff{1}{2}$.
  \item On considère le programme suivant.
\begin{python}
from math import exp

def dichotomie(n):
    a = 1
    b = 2
    for i in range(n):
        m = (a + b) / 2
        if exp(m) > 3:
            b = m
        else:
            a = m
    return a, b
\end{python}
        Faire tourner l'appel \texttt{dichotomie(4)} à la main, en notant à chaque
        passage les valeurs de $m$, $a$ et $b$.
  \item Que renvoie \texttt{dichotomie(4)} ? Interpréter.
  \item Combien de passages faut-il pour un encadrement de $\alpha$ d'amplitude au plus
        $10^{-3}$ ?
\end{enumerate}

\exo[Compléter un programme]
Soit la fonction $f$ définie sur $\R$ par $f(x)=x^{3}+x-3$ ; on admet que l'équation
$f(x)=0$ a une unique solution $\alpha$ dans $\intff{1}{2}$.
\begin{enumerate}
  \item Compléter le programme pour que \texttt{dichotomie(a, b, e)} renvoie un
        encadrement de $\alpha$ d'amplitude au plus \texttt{e}.
\begin{python}
def f(x):
    return x**3 + x - 3

def dichotomie(a, b, e):
    while .......... :
        m = (a + b) / 2
        if f(a) * f(m) <= 0:
            b = ..........
        else:
            a = ..........
    return a, b
\end{python}
  \item Que renvoie \texttt{dichotomie(1, 2, 0.1)} ? Détailler les étapes.
\end{enumerate}

\exo[Balayage ou dichotomie ?]
Soit la fonction $g$ définie sur $\R$ par $g(x)=x^{3}+3x-1$.
\begin{enumerate}
  \item Montrer que l'équation $g(x)=0$ admet une unique solution $\alpha$ dans
        $\intff{0}{1}$.
  \item Par balayage au pas de $0{,}1$ (tableau de valeurs de la calculatrice), donner un
        encadrement de $\alpha$ d'amplitude $0{,}1$.
  \item Un balayage au pas de $0{,}01$ sur $\intff{0}{1}$ demande jusqu'à $100$ calculs
        d'images. Combien de passages de la dichotomie, en partant de $\intff{0}{1}$,
        suffisent pour une amplitude au plus $0{,}01$ ?
\end{enumerate}

\partiefichesansnum{Exercices type bac}

\exo[D'après bac 2026, Métropole]
Soit $f$ la fonction dérivable sur $\R$ définie par $f(x)=(3x-1)\e^{-x}+2$. On note $\Cf$
sa courbe représentative dans un repère orthonormé.
\begin{enumerate}
  \item Déterminer la limite de la fonction $f$ en $-\infty$.
  \item Déterminer la limite de la fonction $f$ en $+\infty$. Interpréter graphiquement
        ce résultat.
  \item On note $f'$ la fonction dérivée de $f$.
  \begin{enumerate}[label=\alph*)]
    \item Vérifier que, pour tout nombre réel $x$, $f'(x)=(4-3x)\e^{-x}$.
    \item Étudier le signe de la dérivée et en déduire le tableau de variations de la
          fonction $f$ sur $\R$.
  \end{enumerate}
  \item
  \begin{enumerate}[label=\alph*)]
    \item Démontrer que l'équation $f(x)=0$ admet une unique solution sur $\R$. On
          notera $\alpha$ cette solution.
    \item Donner une valeur approchée de $\alpha$ à $10^{-2}$ près.
    \item En déduire le tableau de signes de la fonction $f$ sur $\R$.
  \end{enumerate}
  \item On admet que la droite d'équation $y=4x+1$ est la tangente à la courbe $\Cf$ au
        point d'abscisse $0$. Vérifier que, pour tout nombre réel $x$,
        $(3x-1)\e^{-x}+2\leq 4x+1$.
\end{enumerate}

\exo[D'après bac 2025, Amérique du Nord]
On considère la fonction $f$ définie sur $\R$ par $f(x)=x\e^{-x}+2x-1$. On admet que la
fonction $f$ est deux fois dérivable sur $\R$. On appelle $\Cf$ sa courbe représentative
dans un repère orthogonal du plan.
\begin{enumerate}
  \item Déterminer les limites de la fonction $f$ en $-\infty$ et en $+\infty$.
  \item Pour tout nombre réel $x$, calculer $f'(x)$.
  \item Montrer que pour tout nombre réel $x$ : $f''(x)=(x-2)\e^{-x}$.
  \item Étudier la convexité de la fonction $f$.
  \item Étudier les variations de la fonction $f'$ sur $\R$, puis dresser son tableau de
        variations en y faisant apparaître la valeur exacte de l'extremum. Les limites
        de la fonction $f'$ aux bornes de l'intervalle de définition ne sont pas
        attendues.
  \item En déduire le signe de la fonction $f'$ sur $\R$, puis justifier que la fonction
        $f$ est strictement croissante sur $\R$.
  \item Justifier qu'il existe un unique nombre réel $\alpha$ tel que $f(\alpha)=0$.
        Donner un encadrement de $\alpha$, au centième près.
  \item On considère la droite $\Delta$ d'équation $y=2x-1$. Étudier la position
        relative de la courbe $\Cf$ par rapport à la droite $\Delta$.
\end{enumerate}

\exo[D'après bac 2025, Asie]
On considère la fonction $f$ définie sur l'intervalle $\intoo{0}{+\infty}$ par
\begin{center}
$f(x)=\dfrac{\e^{\sqrt{x}}}{2\sqrt{x}}$
\end{center}
et on appelle $\Cf$ sa courbe représentative dans un repère orthonormé.
\begin{enumerate}
  \item On définit la fonction $g$ sur l'intervalle $\intoo{0}{+\infty}$ par
        $g(x)=\e^{\sqrt{x}}$.
  \begin{enumerate}[label=\alph*)]
    \item Montrer que $g'(x)=f(x)$ pour tout $x$ de l'intervalle $\intoo{0}{+\infty}$.
    \item Pour tout nombre réel $x$ de l'intervalle $\intoo{0}{+\infty}$, calculer $f'(x)$
          et montrer que $f'(x)=\dfrac{\e^{\sqrt{x}}\left(\sqrt{x}-1\right)}{4x\sqrt{x}}$.
  \end{enumerate}
  \item
  \begin{enumerate}[label=\alph*)]
    \item Déterminer la limite de la fonction $f$ en $0$.
    \item Interpréter graphiquement ce résultat.
  \end{enumerate}
  \item
  \begin{enumerate}[label=\alph*)]
    \item Déterminer la limite de la fonction $f$ en $+\infty$.
    \item Étudier le sens de variation de la fonction $f$ sur $\intoo{0}{+\infty}$.
          Dresser le tableau de variations de la fonction $f$ en y faisant figurer les
          limites aux bornes de l'intervalle de définition.
    \item Montrer que l'équation $f(x)=2$ admet une unique solution sur l'intervalle
          $\intfo{1}{+\infty}$ et donner une valeur approchée à $10^{-1}$ près de cette
          solution.
  \end{enumerate}
  \item On admet que la fonction $f$ est deux fois dérivable sur l'intervalle
        $\intoo{0}{+\infty}$ et que
        $f''(x)=\dfrac{\e^{\sqrt{x}}\left(x-3\sqrt{x}+3\right)}{8x^{2}\sqrt{x}}$.
  \begin{enumerate}[label=\alph*)]
    \item En posant $X=\sqrt{x}$, montrer que $x-3\sqrt{x}+3>0$ pour tout nombre réel $x$
          de l'intervalle $\intoo{0}{+\infty}$.
    \item Étudier la convexité de la fonction $f$ sur l'intervalle $\intoo{0}{+\infty}$.
  \end{enumerate}
\end{enumerate}

\partiefichesansnum{Entraînement au calcul de dérivées}

\textit{Dans ces exercices, chaque fonction est dérivable sur l'intervalle indiqué.
Écrire la dérivée sous la forme la plus simple, factorisée quand c'est possible.}

\exo[Polynômes, inverses, racines]
Calculer la dérivée de chaque fonction.
\begin{enumerate}[label=\alph*)]
  \item $f$ sur $\R$ : $f(x)=-3x^{2}+8x-5$.
  \item $g$ sur $\intoo{0}{+\infty}$ : $g(x)=\dfrac{x^{3}}{6}-2\sqrt{x}$.
  \item $h$ sur $\R$ : $h(x)=\dfrac{x^{4}-6x^{2}+5x-1}{2}$.
  \item $k$ sur $\R$ : $k(x)=(4x-1)^{2}$.
  \item $m$ sur $\intoo{0}{+\infty}$ : $m(x)=-\dfrac{3}{x^{2}}$.
  \item $p$ sur $\intoo{0}{+\infty}$ : $p(x)=\dfrac{3}{\sqrt{x}}$.
  \item $r$ sur $\intoo{-\infty}{3}$ : $r(x)=5x-2+\dfrac{1}{3-x}$.
  \item $s$ sur $\intoo{0}{+\infty}$ : $s(x)=\dfrac{2}{x}-\dfrac{x^{2}}{4}$.
\end{enumerate}

\exo[Quotients]
Calculer la dérivée de chaque fonction.
\begin{enumerate}[label=\alph*)]
  \item $f$ sur $\intoo{\frac{5}{4}}{+\infty}$ : $f(x)=\dfrac{2-3x}{4x-5}$.
  \item $g$ sur $\intoo{0}{+\infty}$ : $g(x)=\dfrac{4x-3}{x^{3}}$.
  \item $h$ sur $\intoo{\frac{3}{2}}{+\infty}$ : $h(x)=\dfrac{x^{2}-5x+1}{2x-3}$.
  \item $k$ sur $\intoo{-2}{+\infty}$ : $k(x)=\dfrac{3x}{2+x}$.
  \item $m$ sur $\intoo{2}{+\infty}$ : $m(x)=\dfrac{x^{2}+x+4}{x-2}$.
  \item $p$ sur $\intoo{3}{+\infty}$ : $p(x)=2+\dfrac{4x}{x-3}$.
  \item $r$ sur $\R$ : $r(x)=\dfrac{x^{2}+1}{x^{2}+3}$.
  \item $s$ sur $\R$ : $s(x)=\dfrac{x^{2}-x+2}{2x^{2}+1}$.
  \item $w$ sur $\intoo{5}{+\infty}$ : $w(x)=\dfrac{-x^{2}+4x+1}{x^{2}-4x-5}$.
\end{enumerate}

\exo[Avec $\e^{x}$]
Calculer la dérivée de chaque fonction.
\begin{enumerate}[label=\alph*)]
  \item $f$ sur $\R$ : $f(x)=\e^{x}-2x^{2}$.
  \item $g$ sur $\R$ : $g(x)=(x+3)\,\e^{x}$.
  \item $h$ sur $\R$ : $h(x)=4x-3\,\e^{x}$.
  \item $k$ sur $\R$ : $k(x)=\left(1-x^{2}\right)\e^{x}$.
  \item $m$ sur $\R$ : $m(x)=\left(x^{2}-2x+3\right)\e^{x}$.
  \item $p$ sur $\R$ : $p(x)=x^{2}\,\e^{x}$.
  \item $r$ sur $\intoo{-3}{+\infty}$ : $r(x)=\dfrac{\e^{x}}{x+3}$.
  \item $s$ sur $\R$ : $s(x)=\dfrac{x-1}{\e^{x}}$.
  \item $w$ sur $\R$ : $w(x)=\dfrac{\e^{x}}{\e^{x}+2}$.
  \item $z$ sur $\R$ : $z(x)=\dfrac{\e^{x}+3}{\e^{x}}$.
\end{enumerate}

\exo[Avec $\e^{u}$]
Calculer la dérivée de chaque fonction.
\begin{enumerate}[label=\alph*)]
  \item $f$ sur $\R$ : $f(x)=\e^{2x}-5$.
  \item $g$ sur $\R$ : $g(x)=8\,\e^{-0{,}25x}$.
  \item $h$ sur $\R$ : $h(x)=x\,\e^{-2x}$.
  \item $k$ sur $\R$ : $k(x)=\e^{-x^{2}+3x}$.
  \item $m$ sur $\R$ : $m(x)=(3x+1)\,\e^{-0{,}5x}$.
  \item $p$ sur $\R$ : $p(x)=(4-0{,}5x)\,\e^{2x}$.
  \item $r$ sur $\R$ : $r(x)=2x\,\e^{-x+3}$.
  \item $s$ sur $\R$ : $s(x)=5\,\e^{2-x^{2}}$.
  \item $w$ sur $\R$ : $w(x)=\e^{\frac{3-x}{4}}$.
  \item $z$ sur $\intoo{2}{+\infty}$ : $z(x)=\e^{\frac{x+1}{x-2}}$.
  \item $\varphi$ sur $\R$ : $\varphi(x)=\dfrac{1}{\sqrt{2\pi}}\,\e^{-\frac{x^{2}}{2}}$.
\end{enumerate}

\exo[Puissances et racines composées]
Calculer la dérivée de chaque fonction.
\begin{enumerate}[label=\alph*)]
  \item $f$ sur $\R$ : $f(x)=\left(x^{2}-3x+5\right)^{4}$.
  \item $g$ sur $\intoo{-\frac{2}{3}}{+\infty}$ : $g(x)=\dfrac{1}{(3x+2)^{3}}$.
  \item $h$ sur $\intoo{\frac{1}{5}}{+\infty}$ : $h(x)=\sqrt{5x-1}$.
  \item $k$ sur $\R$ : $k(x)=\sqrt{x^{2}-4x+9}$.
  \item $m$ sur $\R$ : $m(x)=(2-x)^{5}$.
  \item $p$ sur $\R$ : $p(x)=\sqrt{\e^{x}+1}$.
  \item $r$ sur $\R$ : $r(x)=\left(\e^{x}-x\right)^{2}$.
  \item $s$ sur $\intoo{-\frac{3}{2}}{+\infty}$ : $s(x)=\dfrac{4}{\sqrt{2x+3}}$.
\end{enumerate}

\exo[Produits et quotients de composées]
Calculer la dérivée de chaque fonction.
\begin{enumerate}[label=\alph*)]
  \item $f$ sur $\intoo{0}{+\infty}$ : $f(x)=x\sqrt{2x+1}$.
  \item $g$ sur $\intoo{0}{+\infty}$ : $g(x)=\dfrac{x\sqrt{x}}{x+2}$.
  \item $h$ sur $\R$ : $h(x)=(x+1)^{2}\,\e^{-x}$.
  \item $k$ sur $\R$ : $k(x)=\dfrac{\e^{2x}}{x^{2}+1}$.
  \item $m$ sur $\R$ : $m(x)=\dfrac{3}{1+2\,\e^{-x}}$.
  \item $p$ sur $\R$ : $p(x)=(2x-1)^{3}\,\e^{x}$.
  \item $r$ sur $\R$ : $r(x)=x^{2}\,\e^{-x^{2}}$.
  \item $s$ sur $\intoo{0}{+\infty}$ : $s(x)=\sqrt{x}\,\e^{-x}$.
\end{enumerate}

\exo[Dériver pour étudier les variations]
Étudier les variations de chaque fonction sur l'intervalle indiqué.
\begin{enumerate}[label=\alph*)]
  \item $f$ sur $\R$ : $f(x)=3\,\e^{-4x}$.
  \item $g$ sur $\R$ : $g(x)=50\,\e^{-0{,}2x+1}$.
  \item $h$ sur $\intff{0}{5}$ : $h(x)=0{,}05\,\e^{2x}+3x$.
  \item $k$ sur $\intff{0}{3}$ : $k(x)=(3-x)\,\e^{\frac{x}{3}}$.
  \item $m$ sur $\intff{-4}{8}$ : $m(x)=\dfrac{20}{1+\e^{-0{,}5x}}$.
  \item $p$ sur $\intoo{-1}{+\infty}$ : $p(x)=\e^{\frac{x^{2}}{x+1}}$.
  \item $r$ sur $\R$ : $r(x)=\left(x^{2}-3\right)\e^{x}$.
\end{enumerate}

\end{multicols}

\end{document}
